\subsection{Rationale}

A locus is strongly differentiated when its allele frequency varies among populations. Two loci show among-population LD only when their frequencies vary across the \emph{same} populations in a consistent direction. Strong differentiation can therefore accumulate across the genome without strong multilocus LD if ancestry has sorted differently among loci. We refer to this as locus-specific ancestry sorting.

\subsection{Units, populations and orientation}

We analysed 165 individuals from 20 hybrid populations. The distance analysis used 20,807 LD-reduced units representing 51,612 ancestry-informative SNPs (parental diagnostic index $>-25$; hybrids-only LD clustering with a decay-relative quality gate $\rho=0.5$ and a 5~cM merge cap). Combining correlated SNPs avoids counting one genomic region repeatedly but removes some short-range LD, so the neighbourhood analysis was repeated with all 51,612 unpruned SNPs. Genotypes were oriented so that dosage counts the \Fa\ allele, identified from the parental reference samples. For each marker we calculated its \Fa-allele frequency in every population and estimated differentiation among populations using the \citet{weir1984estimating} \FST. Using the project-wide sorting criteria (near-fixation $\phi=0.85$, sorting threshold $\tau=0.6$, binomial $\alpha=0.05$), 1,215 units were sorted toward \Fa\ and 365 toward \Fp.

\subsection{Pairwise statistics}

\paragraph{Ancestry-profile similarity.} For each pair of markers $i$ and $j$, ancestry-profile similarity $c_{ij}$ is the correlation of their oriented allele-frequency profiles across the 20 populations. A positive value means that the same populations were enriched for \Fa\ ancestry at both loci. We repeated the calculation after removing each population's genome-wide ancestry, estimated from all other chromosomes.

\paragraph{Among-population LD and its differentiation ceiling.} Following the decomposition of LD in subdivided populations by \citet{ohta1982linkage}, the among-population component of LD between two loci, normalised as a correlation, is
\begin{equation}
r_{ST,ij}=\frac{\operatorname{cov}_k(F_{ik},F_{jk})}{\sqrt{\bar F_i(1-\bar F_i)\,\bar F_j(1-\bar F_j)}}=c_{ij}\sqrt{G_iG_j},
\qquad G_i=\frac{\operatorname{var}_k(F_{ik})}{\bar F_i(1-\bar F_i)},
\label{eq:rst}
\end{equation}
where $G_i$ is the among-population differentiation of locus $i$ (an unweighted, Nei-type measure computed from population frequencies, distinct from the Weir--Cockerham \FST\ used elsewhere). The product $G_iG_j$ is the \emph{ceiling}: the maximum $r_{ST,ij}^2$, reached if the two loci differentiate the same populations perfectly ($|c_{ij}|=1$). We summarised the \emph{realised fraction} of the ceiling as $\sum r_{ST}^2/\sum G_iG_j$.

\paragraph{Permutation baseline.} With only 20 populations, unrelated profiles still have a non-zero expected squared correlation ($1/19$). We estimated the baseline for each distance class by independently permuting population labels for every unit 500 times. This preserved each unit's allele-frequency profile and differentiation while removing correspondence between loci.

\paragraph{Within-population LD.} We correlated oriented genotypes after centring them within populations and removing variation associated with each individual's genome-wide ancestry, estimated from all other chromosomes. Missing genotypes (2.7\%) were set to the population mean. Because sorted loci are near-fixed in many populations, comparisons of sorted and unsorted loci included only populations in which both markers segregated ($r_w^{\mathrm{poly}}$).

\subsection{Distance analysis}

We evaluated all 9.6 million within-chromosome pairs of LD-reduced units and all 206.8 million cross-chromosome pairs. Pairs were grouped by genetic distance, with physical distance used as a sensitivity analysis; cross-chromosome pairs formed an ``unlinked'' class. Confidence intervals came from 2,000 chromosome-block bootstrap replicates.

\subsection{Neighbourhoods of sorted loci}

If sorting involved long haplotypes, markers near sorted loci should show more similar ancestry patterns and stronger within-population LD than markers near comparable unsorted loci. We matched each sorted locus to up to five unsorted controls based on \FST, local recombination rate, parental allele-frequency difference and local density of LD units (units within $\pm100$~kb). A 1-SD caliper retained 1,552 of 1,577 eligible sorted loci and achieved good balance on all covariates (absolute standardised mean differences $\leq0.04$); narrower calipers excluded many large sorted blocks without improving balance. The 25 unmatched loci are listed separately (Table~\ref{tab:unmatched}).

For each focal locus, we averaged ancestry-profile similarity and $r_w^{\mathrm{poly}}$ with neighbouring markers up to 2~Mb away. We did this both for other LD-reduced units and for all unpruned SNPs, including SNPs in the focal locus's own LD cluster. We also compared the number of SNPs and physical span of the focal LD units. Each sorted locus was compared with the mean of its own controls, and 2,000 chromosome-block bootstrap replicates kept these matched sets intact.
